\chapter{Case Study: The Environmental Application}
\label{ch:case-study}

This chapter applies the entire theoretical and mechanical apparatus
developed in Parts~I and~II to the repository's environmental application
example --- the closest the documentation comes to a complete, realistic
environmental computational science workflow. Environmental computational
science is the use case of the whole book: the Gulf of Naples stations of the
lectures are a small model of the forecasting chain described here.

\section{Source and Scope}

\begin{center}
\textit{Source directory:} \texttt{examples/envapp}
\end{center}

\begin{lstlisting}[style=bash, caption={Layout of the environmental application example.}]
examples/envapp/
  app.py
  README.md
  models/
  models/wrf/
  models/wrf/scripts/
  models/wrf/utils/
  archive2plot/
  publishWrf/
\end{lstlisting}

\begin{cautionbox}
The repository describes this material explicitly as a \emph{mockup/sanitized
demonstration} of a larger, real environmental forecasting system. It is not
the production code running at the University of Naples ``Parthenope''
Center for Monitoring and Modeling Marine and Atmosphere applications;
treat it as a structurally faithful template, not as a deployable forecasting
system.
\end{cautionbox}

\section{Scientific Purpose}

The environmental application examples are research-oriented demonstration
material inspired by real forecasting workflows involving three families of
numerical models:

\begin{itemize}[leftmargin=1.4em]
  \item \textbf{WRF} --- the Weather Research and Forecasting model, an
    atmospheric numerical weather prediction system;
  \item \textbf{ROMS} --- the Regional Ocean Modeling System, a
    free-surface, terrain-following ocean circulation model;
  \item \textbf{WaComM++} --- a marine pollution transport-and-diffusion
    model used for contaminant and Lagrangian particle tracking.
\end{itemize}

\section{Concepts Exercised}

The example exercises, in combination, every architectural concept developed
in this book:

\begin{itemize}[leftmargin=1.4em]
  \item wrapping external scientific model executables
    (\Cref{ch:external-software});
  \item preparing model input files;
  \item executing model stages, potentially across \texttt{BATCH},
    \texttt{SLURM}, and \texttt{DOCKER} backends (\Cref{ch:backends});
  \item post-processing outputs;
  \item rendering archived products to plots;
  \item publishing products, naturally suited to a \texttt{WEB} task
    (\Cref{ch:web-tasks}).
\end{itemize}

\section{The WRF Pipeline as a Canonical Numerical-Weather Workflow}

The \texttt{models/wrf/scripts/} directory includes shell-oriented wrappers
representing the standard structure of a numerical weather workflow:

\begin{center}
\begin{tabular}{@{}ll@{}}
\toprule
\textbf{Script family} & \textbf{Role} \\
\midrule
\texttt{geogrid}  & Define and prepare the model's geographical domain. \\
\texttt{ungrib}   & Decode gridded meteorological input (e.g.\ GRIB) data. \\
\texttt{metgrid}  & Horizontally interpolate meteorological data onto the domain. \\
\texttt{real}     & Vertically interpolate and generate initial/boundary conditions. \\
\texttt{wrf}      & Execute the core atmospheric simulation. \\
\texttt{publishWrfOutput} & Export and disseminate simulation products. \\
namelist generators & Produce the model's run-time configuration files. \\
\bottomrule
\end{tabular}
\end{center}

This script family corresponds to the canonical six-stage structure of a
WRF-based forecast:

\begin{enumerate}[leftmargin=1.6em]
  \item configure domain and static geographical data;
  \item prepare meteorological input;
  \item generate model-ready files;
  \item run initialization;
  \item run simulation;
  \item post-process and publish outputs.
\end{enumerate}

\begin{figure}[h]
\centering
\begin{tikzpicture}[
  node distance=8mm,
  task/.style={rectangle, rounded corners, draw=dagonblue, very thick,
    fill=blue!4, minimum width=20mm, minimum height=10mm, font=\sffamily\scriptsize, align=center},
  ext/.style={rectangle, rounded corners, draw=dagonaccent, very thick,
    fill=teal!6, minimum width=20mm, minimum height=10mm, font=\sffamily\scriptsize, align=center},
  arr/.style={-{Stealth[length=2.2mm]}, thick, dagonblue}
]
  \node[task] (geo) {geogrid};
  \node[task, right=14mm of geo] (ung) {ungrib};
  \node[task, right=14mm of ung] (met) {metgrid};
  \node[task, right=14mm of met] (real) {real};
  \node[task, right=14mm of real] (wrf) {wrf\\ \footnotesize (\texttt{SLURM})};
  \node[task, right=14mm of wrf] (a2p) {archive2plot};
  \node[ext, right=14mm of a2p] (pub) {publishWrf\\ \footnotesize (\texttt{WEB})};
  \draw[arr] (geo) -- (met);
  \draw[arr] (ung) -- (met);
  \draw[arr] (met) -- (real);
  \draw[arr] (real) -- (wrf);
  \draw[arr] (wrf) -- (a2p);
  \draw[arr] (a2p) -- (pub);
\end{tikzpicture}
\caption{A DAGonStar-style DAG induced by the canonical WRF pipeline.
\texttt{geogrid} and \texttt{ungrib} are independent producers consumed by
\texttt{metgrid}, after which the chain becomes linear through simulation,
rendering, and publication. Backend annotations indicate a representative
(not prescriptive) execution-environment assignment, following
\Cref{tab:backend-selection}.}
\label{fig:wrf-pipeline}
\end{figure}

\section{Post-Processing: \texttt{archive2plot}}

\texttt{archive2plot} contains a utility for rendering post-processed model
output as image plots:

\begin{lstlisting}[style=bash]
archive2plot/
  archive2plot.py
  Render.py
  maps.json
  archive2plot.conf
  requirements.txt
  (shapefile documentation)
\end{lstlisting}

In DAGonStar terms, \texttt{archive2plot} is naturally either a
\texttt{BATCH} task (if invoked as a standalone Python script with its own
environment already on the host) or a \texttt{NATIVE} task
(\Cref{ch:native-tasks}), if its rendering logic is refactored into an
importable, module-level function consuming a staged NetCDF product and
producing a declared image output.

\section{Prerequisites and Operational Honesty}

The documentation is explicit that these examples may require substantial
external software:

\begin{itemize}[leftmargin=1.4em]
  \item compiled WRF and related preprocessing tools;
  \item geospatial Python libraries;
  \item NetCDF/HDF libraries;
  \item model input datasets;
  \item shapefiles or map assets;
  \item HPC or containerized execution environments.
\end{itemize}

\begin{cautionbox}
These are not generic unit tests. They should be treated as scientific
application \emph{templates} --- structurally complete, but requiring
site-specific software, data, and credentials before they can be executed
end to end. This is consistent with the project's own examples-catalog
guidance (\Cref{ch:examples-catalog}), which distinguishes locally verifiable
examples from those requiring external software or services.
\end{cautionbox}

\section{Adapting the Pattern: A Five-Step Method}

The documentation's own adaptation checklist, already previewed structurally
in \Cref{ch:external-software}'s five-stage pattern, applies directly here:

\begin{enumerate}[leftmargin=1.6em]
  \item identify each model executable as a task;
  \item define all file products exchanged between tasks;
  \item wrap each executable in a reproducible launcher script;
  \item encode cross-task data products with \texttt{workflow://};
  \item verify each model stage independently before running the full
    workflow.
\end{enumerate}

\begin{example}
Applying step~5 concretely to \Cref{fig:wrf-pipeline}: before composing the
full seven-task DAG, a researcher should first verify \texttt{geogrid} alone
(does it produce the expected static geographical files?), then
\texttt{ungrib} alone (does it correctly decode the chosen meteorological
input source?), and only then connect them through \texttt{metgrid} via
\texttt{workflow://} references. This staged verification strategy mirrors
the local-structural-verification pattern already introduced for Slurm and
remote tasks in \Cref{ch:backends}: verify each component in isolation before
trusting the composed graph.
\end{example}

\section{Real-World Anchoring}

The environmental application generalizes patterns from real, funded
deployments that the project's own acknowledgments document
(reproduced in full in the Bibliography):

\begin{itemize}[leftmargin=1.4em]
  \item \textbf{MytilX} --- a research agreement on modelling mytilus
    (mussel) farming at scale, in which DAGonStar orchestrates the
    production workflow delivering study-zone-specific weather, marine,
    pollutant, and farmed-mussel contamination forecasts and predictions.
  \item \textbf{MytilEx} --- a research contract on mussel farming with HPC
    and AI, in which DAGonStar orchestrates the daily production of 168
    forecast hours of weather, marine, pollutant, and contamination
    predictions.
  \item \textbf{ADMIRE} (EuroHPC H2020, Work Package~7) --- DAGonStar
    orchestrates an environmental application delivering on-demand weather,
    marine, and pollutant simulations and forecasts for the Campania Region,
    Italy.
\end{itemize}

These deployments are the operational referent for the synthesis chapter
that follows: real systems in which DAGonStar coordinates exactly the kind of
numerical-model, post-processing, and publication chain sketched in
\Cref{fig:wrf-pipeline}, at production cadence.

\section{Reading the WRF Pipeline as a Layered Dependency Structure}

\Cref{fig:wrf-pipeline} presented the WRF pipeline as a DAG; it is worth
examining more closely \emph{why} the pipeline has the particular dependency
shape it does, since the shape directly reflects the scientific structure
of numerical weather prediction, not an arbitrary engineering choice.

\texttt{geogrid} and \texttt{ungrib} are independent of one another because
they consume entirely disjoint kinds of input: \texttt{geogrid} processes
static geographical data (terrain elevation, land use, soil type) that does
not depend on any particular forecast initialization time, while
\texttt{ungrib} decodes time-varying meteorological fields (temperature,
pressure, wind) from an external data source. Because neither task's output
depends on the other's, the engine is free, per the topological-ordering
discussion of \Cref{sec:topological-ordering}, to execute them concurrently
--- a genuine, not merely theoretical, opportunity for parallelism that a
researcher decomposing this pipeline by hand should explicitly preserve by
\emph{not} introducing an artificial explicit dependency between the two
merely because they happen to be related model-preparation steps.

\texttt{metgrid} depends on \emph{both}, because horizontal interpolation
onto the model's domain genuinely requires both the static geography
\texttt{geogrid} produced and the decoded meteorological fields
\texttt{ungrib} produced; this is a genuine data dependency in the sense of
\Cref{sec:task-vs-data}, not a procedural one, and is correctly expressed,
in a DAGonStar implementation of this pipeline, through two separate
\texttt{workflow://} references inside \texttt{metgrid}'s command rather
than through a single combined reference. From \texttt{metgrid} onward, the
chain becomes strictly linear --- \texttt{real}, then \texttt{wrf}, then
post-processing and publication --- because each subsequent stage
genuinely requires the immediately preceding stage's specific numerical
output and nothing else.

\begin{practicebox}
This analysis generalizes into a useful design heuristic for decomposing any
new, unfamiliar scientific pipeline into a DAGonStar workflow: before
writing any \texttt{DagonTask} code, first identify, for each candidate
stage, the \emph{minimal} set of upstream products it genuinely requires,
and resist the temptation to introduce additional, unnecessary dependency
edges purely because two stages are conceptually related or were
historically run in a fixed order out of habit rather than necessity. Every
unnecessary edge removes a genuine opportunity for the engine to exploit
concurrency, exactly as an artificial \texttt{geogrid}~$\to$~\texttt{ungrib}
edge would needlessly serialize two genuinely independent stages of the
pipeline examined here.
\end{practicebox}

\section{Resource Heterogeneity Within a Single Pipeline}

A second dimension worth making explicit, beyond the dependency structure
already analyzed, is the heterogeneity of \emph{computational resource
requirements} across the WRF pipeline's stages --- a heterogeneity that
directly motivates the backend-selection guidance of
\Cref{tab:backend-selection} and that a researcher decomposing a similar
pipeline should reason about explicitly rather than defaulting every stage
to the same backend out of convenience.

\texttt{geogrid} and \texttt{ungrib} are typically I/O-bound, data-
preparation steps whose resource demands are modest relative to the core
simulation; running them as local \texttt{BATCH} tasks, or lightweight
remote \texttt{BATCH} tasks, is usually entirely adequate. \texttt{wrf}
itself, the core atmospheric simulation, is typically the pipeline's most
computationally demanding stage by a wide margin, routinely requiring
many-core, multi-node parallel execution over a span of hours --- precisely
the profile \texttt{TaskType.SLURM} exists to address, with its explicit
\texttt{ntasks}, \texttt{nodes}, and \texttt{time} resource-request
parameters (\Cref{ch:backends}) mapping directly onto the simulation's
actual computational needs. \texttt{archive2plot}'s rendering step sits
between these two extremes: more demanding than simple data preparation,
since it processes the simulation's full output volume, but typically far
less demanding than the simulation itself, and frequently dependent on a
specific geospatial software stack (the shapefile and mapping libraries
noted in \Cref{ch:case-study}'s file listing) that argues for
\texttt{TaskType.DOCKER} on dependency-management grounds, independent of
its computational intensity.

\begin{practicebox}
The general principle this illustrates: a realistic scientific pipeline's
backend assignment should follow the \emph{resource and dependency profile
of each individual stage}, not a single, pipeline-wide default. The
environmental application example's implicit backend diversity --- light
local or remote tasks for preparation, Slurm for the simulation core, Docker
for a specialized rendering toolchain --- is not an arbitrary illustrative
choice in \Cref{fig:wrf-pipeline}; it is the natural consequence of applying
\Cref{tab:backend-selection}'s selection criteria stage by stage to a
pipeline whose stages genuinely differ in their computational and
dependency-management needs.
\end{practicebox}

\section{From Demonstration Pipeline to Operational Deployment}

\Cref{ch:case-study} was careful to distinguish the sanitized, demonstration
\texttt{examples/envapp} material from the production deployments it
generalizes (MytilX, MytilEx, ADMIRE), and it is worth examining, in more
depth than the original chapter's brief mention, what additional concerns
typically separate a demonstration pipeline from an operational one running
at production cadence --- concerns this book's earlier chapters have each
already introduced individually, but that converge specifically at the
point of operational deployment.

An operational, daily-cadence forecasting chain must additionally address:
\textbf{scheduling} --- something external to DAGonStar itself (a cron job,
a system-level scheduler, or a controller using the asynchronous-launch
pattern of \Cref{ch:meta-async}) must trigger a new workflow run at the
appropriate cadence; \textbf{failure recovery across runs}, not merely
within a single run --- checkpointing (\Cref{ch:lessons-05-07}) addresses
resuming a single failed run, but an operational deployment must also
decide what happens if an entire day's forecast cycle fails outright, a
question outside the scope of any single workflow's own configuration;
\textbf{monitoring and alerting} --- the lifecycle event hooks of
\Cref{ch:meta-async} provide the mechanism, but an operational deployment
must build the actual alerting logic on top of them; and \textbf{data
retention policy} --- \Cref{ch:lessons-05-07}'s checkpoint-persistence
discussion already noted that checkpoint reuse depends on scratch
directories remaining on disk, which has direct, compounding storage-cost
implications once a pipeline runs daily indefinitely rather than once, for
a tutorial exercise.

\begin{cautionbox}
None of these four operational concerns is addressed by DAGonStar's own
documented feature set in isolation; each requires deliberate engineering
on top of the mechanisms this book has developed. A researcher planning to
move from a demonstration pipeline, of the kind examined in this chapter, to
genuine operational deployment should treat this list as a minimum
checklist of additional design work required, not as a list of problems
DAGonStar itself solves automatically.
\end{cautionbox}

\section{Chapter Summary}

The environmental application example is the richest demonstration of
DAGonStar's intended domain: orchestrating heterogeneous model components and
data products across scientific software stacks. This chapter additionally
analyzed the WRF pipeline's dependency structure as a direct reflection of
the underlying numerical weather prediction science rather than an arbitrary
engineering choice, examined the resource heterogeneity across pipeline
stages that motivates differentiated backend selection, and enumerated the
additional operational concerns --- scheduling, cross-run failure recovery,
monitoring, and data retention --- that separate a demonstration pipeline
from a genuine, production-cadence operational deployment of the kind the
MytilX, MytilEx, and ADMIRE acknowledgments describe. \Cref{ch:examples-catalog}
now surveys the remainder of the project's example catalog, and
\Cref{ch:synthesis} closes the book by assembling these pieces into the
shape of a complete, AI-augmented, operational forecasting chain.

\begin{exercise}
Using the dependency-structure analysis of this chapter's first section,
identify one further pair of stages in \Cref{fig:wrf-pipeline}, beyond
\texttt{geogrid}/\texttt{ungrib}, that could in principle run concurrently
if the pipeline were extended to include a second, independent
post-processing path (for instance, a statistical verification step
computed directly from \texttt{wrf}'s output, alongside
\texttt{archive2plot}'s rendering step). Sketch the resulting DAG.
\end{exercise}
